% Bidirectional Higher-Rank Type System
% From Practical Type Inference for Higher-Rank Types (2007)

\documentclass[11pt,a4paper]{article}

\usepackage{amsmath,amssymb,amsthm}
\usepackage{mathpartir}
\usepackage{geometry}

\geometry{margin=1in}

\begin{document}

\title{Bidirectional Type System: Lambda Rules}
\date{}
\maketitle

\section{Overview}

These rules define a \textbf{bidirectional type system} that \textbf{elaborates} (translates) source language terms to System F, a polymorphic lambda calculus. The judgment $t \mapsto e$ indicates that source term $t$ translates to System F term $e$. This elaboration makes implicit type instantiation and generalization explicit.

\section{System F: Target Language}

System F is the polymorphic lambda calculus that serves as the target of elaboration.

\subsection{System F Types}
\[
\begin{array}{rcl}
\sigma & ::= & a \mid \sigma_1 \rightarrow \sigma_2 \mid \forall a.\,\sigma \quad \text{(Polymorphic types)}
\end{array}
\]

\subsection{System F Terms}
\[
\begin{array}{rcl}
e & ::= & x \mid n \mid \lambda x:\sigma.\,e \mid e_1\,e_2 \mid \Lambda a.\,e \mid e\,\sigma \mid \textsf{let}\;x:\sigma = e_1\;\textsf{in}\;e_2
\end{array}
\]

Where:
\begin{itemize}
\item $\lambda x:\sigma.\,e$ is term abstraction (function)
\item $e_1\,e_2$ is term application
\item $\Lambda a.\,e$ is type abstraction (polymorphic introduction)
\item $e\,\sigma$ is type application (polymorphic elimination)
\end{itemize}

\section{Source Language Syntax}

\subsection{Source Terms}
\[
t ::= x \mid n \mid \lambda x.\,t \mid \lambda x\!::\!\sigma.\,t \mid t\,u \mid \textsf{let}\;x = t\;\textsf{in}\;u \mid (t :: \sigma)
\]

\subsection{Types}
\[
\sigma ::= \forall \overline{a}.\,\rho \quad \text{(Polymorphic)}
\]
\[
\rho ::= \tau \mid \sigma_1 \rightarrow \sigma_2 \quad \text{(Rho types)}
\]
\[
\tau ::= a \mid \tau_1 \rightarrow \tau_2 \quad \text{(Monomorphic)}
\]

\subsection{Directions}
\[
\delta ::= \Downarrow \mid \Uparrow \quad \text{(check vs synth)}
\]

\section{Conventions}

\begin{itemize}
\item The direction $\delta$ is placed as a subscript on the turnstile: $\Gamma \vdash_\delta t : \rho \mapsto e$
\item $\vdash_\Downarrow$ denotes \textbf{checking} mode (type is given, check term against it)
\item $\vdash_\Uparrow$ denotes \textbf{synthesis} mode (term is given, synthesize its type)
\item $\mapsto$ denotes \textbf{elaboration}: source term $t$ translates to System F term $e$
\item Superscripts on the turnstile indicate specific judgment variants:
\begin{itemize}
\item $\vdash^{\text{poly}}$ for polymorphic generalization
\item $\vdash^{\text{inst}}$ for instantiation
\end{itemize}
\item $\overline{a}$ denotes a sequence of type variables $a_1, a_2, \ldots, a_n$ (possibly empty)
\end{itemize}

\section{$\Gamma \vdash_\delta t : \rho \mapsto e$}

\begin{mathpar}
\inferrule*[right=INT]
{ }
{\Gamma \vdash_\delta n : \textsf{Int} \mapsto n}

\inferrule*[right=VAR]
{\vdash^{\text{inst}}_\delta \sigma \leq \rho \mapsto f}
{\Gamma, (x:\sigma) \vdash_\delta x : \rho \mapsto f\,x}

\inferrule*[right=ABS1]
{\Gamma, x:\tau \vdash_\Uparrow t : \rho \mapsto e}
{\Gamma \vdash_\Uparrow \lambda x.\,t : \tau \rightarrow \rho \mapsto \lambda x:\tau.\,e}

\inferrule*[right=ABS2]
{\Gamma, x:\sigma_a \vdash^{\text{poly}}_\Downarrow t : \sigma_r \mapsto e}
{\Gamma \vdash_\Downarrow \lambda x.\,t : \sigma_a \rightarrow \sigma_r \mapsto \lambda x:\sigma_a.\,e}

\inferrule*[right=AABS1]
{\Gamma, x:\sigma \vdash_\Uparrow t : \rho \mapsto e}
{\Gamma \vdash_\Uparrow \lambda x\!::\!\sigma.\,t : \sigma \rightarrow \rho \mapsto \lambda x:\sigma.\,e}

\inferrule*[right=AABS2]
{\vdash^{\text{dsk}} \sigma_a \leq \sigma_x \mapsto f \\
 \Gamma, x:\sigma_x \vdash^{\text{poly}}_\Downarrow t : \sigma_r \mapsto e}
{\Gamma \vdash_\Downarrow \lambda x\!::\!\sigma_x.\,t : \sigma_a \rightarrow \sigma_r \mapsto \lambda x:\sigma_a.\,[x \mapsto (f\,x)]e}

\inferrule*[right=APP]
{\Gamma \vdash_\Uparrow t : \sigma \rightarrow \sigma' \mapsto e_1 \\
 \Gamma \vdash^{\text{poly}}_\Downarrow u : \sigma \mapsto e_2 \\
 \vdash^{\text{inst}}_\delta \sigma' \leq \rho \mapsto f}
{\Gamma \vdash_\delta t\,u : \rho \mapsto f\,(e_1\,e_2)}

\inferrule*[right=ANNOT]
{\Gamma \vdash^{\text{poly}}_\Downarrow t : \sigma \mapsto e \\
 \vdash^{\text{inst}}_\delta \sigma \leq \rho \mapsto f}
{\Gamma \vdash_\delta (t :: \sigma) : \rho \mapsto f\,e}

\inferrule*[right=LET]
{\Gamma \vdash^{\text{poly}}_\Uparrow \mu : \sigma \mapsto e_1 \\
 \Gamma, x:\sigma \vdash_\delta t : \rho \mapsto e_2}
{\Gamma \vdash_\delta \textsf{let}\;x = \mu\;\textsf{in}\;t : \rho \mapsto \textsf{let}\;x:\sigma = e_1\;\textsf{in}\;e_2}
\end{mathpar}

\section{$\Gamma \vdash^{\text{poly}}_\delta t : \sigma \mapsto e$}

\begin{mathpar}
\inferrule*[right=GEN1]
{\overline{a} = \text{ftv}(\rho) - \text{ftv}(\Gamma) \\
 \Gamma \vdash_\Uparrow t : \rho \mapsto e}
{\Gamma \vdash^{\text{poly}}_\Uparrow t : \forall\overline{a}.\,\rho \mapsto \Lambda\overline{a}.\,e}

\inferrule*[right=GEN2]
{pr(\sigma) = \forall\overline{a}.\,\rho \mapsto f \\
 \overline{a} \notin \text{ftv}(\Gamma) \\
 \Gamma \vdash_\Downarrow t : \rho \mapsto e}
{\Gamma \vdash^{\text{poly}}_\Downarrow t : \sigma \mapsto f\,(\Lambda\overline{a}.\,e)}
\end{mathpar}

\section{$\vdash^{\text{inst}}_\delta \sigma \leq \rho \mapsto f$}

\begin{mathpar}
\inferrule*[right=INST1]
{ }
{\vdash^{\text{inst}}_\Uparrow \forall\overline{a}.\,\rho \leq [\overline{a \mapsto \tau}]\rho \mapsto \lambda x:\forall\overline{a}.\,\rho.\,x\,\overline{\tau}}

\inferrule*[right=INST2]
{\vdash^{\text{dsk}} \sigma \leq \rho \mapsto f}
{\vdash^{\text{inst}}_\Downarrow \sigma \leq \rho \mapsto f}
\end{mathpar}

\section{$pr(\sigma) = \forall\overline{a}.\,\rho \mapsto f$}

\begin{mathpar}
\inferrule*[right=PRPOLY]
{pr(\rho_1) = \forall\overline{b}.\,\rho_2 \mapsto f \\
 \overline{a} \notin \overline{b}}
{pr(\forall\overline{a}.\,\rho_1) = \forall\overline{a}\overline{b}.\,\rho_2 \mapsto \lambda x:\forall\overline{a}\overline{b}.\,\rho_2.\,\Lambda\overline{a}.\,f\,(x\,\overline{a})}

\inferrule*[right=PRFUN]
{pr(\sigma_2) = \forall\overline{a}.\,\rho_2 \mapsto f \\
 \overline{a} \notin \text{ftv}(\sigma_1)}
{pr(\sigma_1 \rightarrow \sigma_2) = \forall\overline{a}.\,(\sigma_1 \rightarrow \rho_2) \mapsto \lambda x:(\forall\overline{a}.\,\sigma_1 \rightarrow \rho_2).\,\lambda y:\sigma_1.\,f\,(\Lambda\overline{a}.\,x\,\overline{a}\,y)}

\inferrule*[right=PRMONO]
{ }
{pr(\tau) = \tau \mapsto \lambda x:\tau.\,x}
\end{mathpar}

\section{$\vdash^{\text{dsk}} \sigma \leq \sigma' \mapsto f$}

\begin{mathpar}
\inferrule*[right=DEEP-SKOL]
{pr(\sigma_2) = \forall\overline{a}.\,\rho \mapsto f_1 \\
 \overline{a} \notin \text{fv}(\sigma_1) \\
 \vdash^{\text{dsk}*} \sigma_1 \leq \rho \mapsto f_2}
{\vdash^{\text{dsk}} \sigma_1 \leq \sigma_2 \mapsto \lambda x:\sigma_1.\,f_1\,(\Lambda\overline{a}.\,f_2\,x)}
\end{mathpar}

\section{$\vdash^{\text{dsk}*} \sigma \leq \rho \mapsto f$}

\begin{mathpar}
\inferrule*[right=SPEC]
{\vdash^{\text{dsk}*} [\overline{a \mapsto \tau}]\rho_1 \leq \rho_2 \mapsto f}
{\vdash^{\text{dsk}*} \forall\overline{a}.\,\rho_1 \leq \rho_2 \mapsto \lambda x:\forall\overline{a}.\,\rho_1.\,f\,(x\,\overline{\tau})}

\inferrule*[right=FUN]
{\vdash^{\text{dsk}} \sigma_3 \leq \sigma_1 \mapsto f_1 \\
 \vdash^{\text{dsk}*} \sigma_2 \leq \sigma_4 \mapsto f_2}
{\vdash^{\text{dsk}*} (\sigma_1 \rightarrow \sigma_2) \leq (\sigma_3 \rightarrow \sigma_4) \mapsto \lambda x:\sigma_1 \rightarrow \sigma_2.\,\lambda y:\sigma_3.\,f_2\,(x\,(f_1\,y))}

\inferrule*[right=MONO]
{ }
{\vdash^{\text{dsk}*} \tau \leq \tau \mapsto \lambda x:\tau.\,x}
\end{mathpar}

\end{document}
